\section{Related Work}
\label{sec:related}

\paragraph{Grammar-based SQL generation.} Closest to us is SQLBull, which starts
generation from grammar rules a tester marks as feature-rich, then builds a path
back to the root so that every query contains the rule; it found 63 previously
unknown defects in five engines, all crashes, internal errors or assertion
failures~\cite{liang2026sqlbull}. SQUIRREL keeps mutated SQL valid and guides the search with coverage
feedback~\cite{zhong2020squirrel}. We
share their goal and their kind of oracle, but we choose where to aim differently.
SQLBull aims at deep grammar, a property of the language. We aim at what a
specific incumbent fuzzer cannot express, a property of its source, and we take
the queries from another engine's bug reports rather than from a grammar walk.

\paragraph{Oracles for database engines.} Ternary logic partitioning splits a
query into three partitions whose union must equal the
original~\cite{rigger2020tlp}. Non-optimizing reference engine construction
compares an optimized query with a rewrite the optimizer cannot
touch~\cite{rigger2020norec}. Both find wrong answers that leave the process
running, which a crash oracle cannot see. We use the weaker oracle on purpose:
it needs neither a reference engine nor an equivalent rewrite, so it can judge
\sql{lpad}, which SQLite does not have.

\paragraph{Keeping generated input valid.} Grammar-based whitebox fuzzing lifts
symbolic constraints to tokens so that solved inputs parse~\cite{godefroid2008grammar}.
Fuzzilli mutates an intermediate language in which every program stays
valid~\cite{gross2023fuzzilli}. Repair-driven greybox fuzzing mutates freely and
repairs the result back to valid syntax~\cite{bendrissou2026repair}. Directed
grammar-based generation learns which input structures reach a testing goal
from feedback~\cite{kirschner2025directed}. FuzzNG fuzzes Linux system calls
without hand-written descriptions, and found defects in components that Syzkaller
had covered for years~\cite{bulekov2023fuzzng}. We solve the validity problem
more crudely: we keep a mined query's structure fixed and change only its
literals, and we learn nothing from feedback. FuzzNG's lesson, that a
long-running incumbent fuzzer still leaves reachable defects, holds here too; we
differ in locating those defects by reading the incumbent's source rather than
by reshaping its input space.

\paragraph{Recurring defects.} VulPA writes a known vulnerability as a pattern
and finds other code that breaks it with multi-object typestate
analysis~\cite{cao2025vulpa}. FirmRec builds signatures of known
vulnerabilities from how they are exploited and finds them again in firmware
images~\cite{xiao2024firmrec}. A recent study measures the errors of
recurring-vulnerability detectors on vulnerabilities inherited through reused
code~\cite{cao2025antman}. Another shows that code which looks
like a known vulnerability but is already patched causes many false
alarms~\cite{tan2025fvf}. We share the question these works ask: does a known
defect have siblings? But they answer it statically, within one code base or
across copied code. We answer it dynamically, across two engines that share no
code. Our duplicate check before filing plays the role that a patched-code
filter plays for them, but each of our answers is a statement that either ends
the process or does not, so there is no false alarm to filter.

\paragraph{Learning from patches.} Mystique ports a vulnerability patch from one
branch to another with a language model~\cite{wu2025mystique}, and VicDiff
locates the commit that introduced a vulnerability by comparing patching
patterns~\cite{guo2026vicdiff}. Both reuse what a patch says about a defect. We
reuse what a fix's regression test says about where defects live, and move a
failing input rather than a fix.
